%  The dissertation abstract can only be 500 words.

Radiation monitoring systems for nuclear environments require coordinated advances across functional materials, detection devices, and radiation-tolerant electronic instrumentation. This dissertation investigates halide perovskite materials, fiber-optic thermal sensors, and FPGA-based digital electronics as complementary technologies for advanced radiation detection and monitoring, organized around a materials-to-systems hierarchy.

Photovoltaic-architecture Cs$_{0.1}$FA$_{0.9}$PbI$_3$ perovskite solar cells were characterized under X-ray, gamma, and neutron irradiation. Under $^{137}$Cs gamma irradiation at 44.16~rad\,hr$^{-1}$, an induced current density of 2.19~nA\,cm$^{-2}$ was measured with a confirmed step response upon source insertion and withdrawal. Extended $^{60}$Co irradiation at 355.61~rad\,hr$^{-1}$ yielded 3.56~nA\,cm$^{-2}$, demonstrating that the current-mode detection approach scales across nearly an order of magnitude in dose rate. A linear degradation rate of approximately 0.74~nA\,rad$^{-1}$ was observed during six-day continuous irradiation, attributable to cumulative defect creation and substrate darkening.

In parallel, perovskite semiconductor radiation detectors were developed for direct-conversion operation. A two-dimensional perovskite X-ray detector incorporating 5F-PEAI surface passivation was characterized, addressing the voltage-induced instability problem caused by ion migration under applied bias. A MAPbI$_3$--hBN composite detector was demonstrated for X-ray detection, and a solution-grown FAPbBr$_3$ single crystal achieved proof-of-concept gamma spectroscopy. For fast neutron detection, the hydrogen-rich MHyPbCl$_3$ perovskite was synthesized and tested at the Ohio State University Research Reactor Fast Beam Facility, exhibiting a resistivity of $4.43 \times 10^{11}$~$\Omega$\,cm and an X-ray sensitivity of 61.74~$\mu$C\,Gy$_{\mathrm{air}}^{-1}$\,cm$^{-2}$. Distinct radiation-induced pulses were recorded from the reactor fast neutron beam, and gamma discrimination analysis using beam filtering confirmed that the observed signals were predominantly attributable to neutron interactions.

As an alternative sensing pathway, thermal radiation sensing was investigated through simulation and experimental testing of an Optical Fiber-Based Gamma Thermometer (OFBGT) and a proposed Optical Fiber-Based Neutron Dosimeter (OFBND). Analytical thermal models predicted temperature differentials of 27--30$\,^{\circ}$C for the OFBGT and 155--159$\,^{\circ}$C for the OFBND. Experimental testing of the OFBGT yielded a measured differential of approximately 17$\,^{\circ}$C, with the discrepancy attributed to the omission of radiative heat transfer in the original simulations.

For digital instrumentation, radiation-tolerant FPGA-based architectures were evaluated through a Flip-Flop Shift Register (FFSR) with dual modular redundancy and a fixed-point thermal control benchmark, both tested under fast neutron irradiation. During two-day testing of five Spartan-7 FPGAs, configuration memory single-event upsets were observed at moderate fluence levels ($\sim$$10^9$~n/cm$^2$), while higher fluences ($\sim$$10^{11}$~n/cm$^2$) produced latch-up events, a multi-SEU cluster, and measurable total ionizing dose degradation. These results demonstrate that configuration memory scrubbing and defense-in-depth strategies are essential for SRAM-based FPGAs in sustained radiation environments.

Together, these investigations span the functional layers of a radiation monitoring system---from radiation-responsive materials and detection devices through alternative thermal sensors to the radiation-tolerant digital electronics required to process their outputs.